\documentclass[10pt,a4paper]{amsart}
\usepackage[margin=19mm]{geometry}
\usepackage{amsmath,amssymb,mathtools}
\usepackage[scr=rsfs,cal=euler]{mathalfa}
\usepackage{xfrac}
\usepackage[unicode,hidelinks,bookmarks=false]{hyperref}
\newcommand{\ie}{i.e.\ }
\newcommand{\cf}{cf.\ }
\newcommand{\resp}{respectively\ }
\newcommand{\Subsection}{\subsection}
\providecommand{\nobreakdash}{\nobreak\hbox{-}}
\setlength{\emergencystretch}{2em}
\multlinegap0pt
\usepackage{amssymb}
\usepackage[shortlabels]{enumitem}
\newtheorem{claim*}{Claim}
\newcommand{\C}{\mathbb{C}}
\newcommand{\ns}{\mathrm{ns}}
\title{Optimal embedding dimension in the Nash--Tognoli theorem}
\author{Juliusz Banecki}
\date{Source: arXiv:2605.22928v1 (CC BY 4.0)}
\begin{document}
\maketitle

\noindent\emph{Source and licence.} Optimal embedding dimension in the Nash--Tognoli theorem. Version arXiv:2605.22928v1 (CC BY 4.0), distributed by arXiv under the Creative Commons Attribution 4.0 International licence, http://creativecommons.org/licenses/by/4.0/, as recorded in the arXiv licence metadata for this exact version and shown on the abstract page https://arxiv.org/abs/2605.22928v1. Full licence notice: \texttt{licenses/arxiv-2605.22928-CC-BY-4.0.txt}.\par\smallskip

\noindent\textit{Editorial scope.} Counted unit: the article's claim* claim* (source0.tex lines 277--279). The article's complete original proof of it is delivered with it (source0.tex lines 280--294, one \texttt{proof} environment of 15 lines). No mathematics is written, repaired or re-derived: the statement and the proof are the article's own lines, in the article's own notation, and no retained line is substituted or re-flowed.  No further source-local context is needed: every cross-reference of every retained line resolves inside the two delivered spans.  Nothing was omitted from any retained span, so the package carries no omission record.  No retained line contains a citation, so the wrapper carries no bibliography and no bibliography entry.  No retained line contains a figure environment, an image-inclusion command or drawing code, so no image is delivered and the package declares no asset dependency.  Editorial formatting only, in addition: an AMS \texttt{amsart} preamble that restates the article's own declarations and macros used by the retained lines, namely the environments \texttt{claim*} on separate counters, together with the standard packages those lines need.  The retained environments print in this document as Claim* 1 in order of appearance, whereas the article prints the same items with the numbering of its own full document.  The complete native source of the article as distributed in the arXiv e-print for this exact version is bundled unchanged as \texttt{sources/201167/source0.tex}.
\begin{claim*}
We have $\mathfrak mB=\mathfrak n$.
\end{claim*}

\begin{proof}
Obviously $\mathfrak mB\subset\mathfrak n$, so it suffices to prove the other inclusion.

The map $f$ is an immersion at $x$, so its derivative is of rank $\dim X$. Hence, the derivative of $f_\C$ at $x$ considered as a linear map over $\C$ is also of rank $\dim X_\C=\dim X$, where by $\dim X_\C$ we mean its complex dimension. It follows that it induces an isomorphism between the tangent spaces to $X_\C$ at $x$ and to $Y_\C$ at $y$. From this we deduce that
\begin{equation*}
    \mathfrak nB_{\mathfrak n}=\mathfrak m B_{\mathfrak n}+\mathfrak n^2B_{\mathfrak n}.
\end{equation*}
Using Nakayama's lemma we find that in fact $\mathfrak nB_{\mathfrak n}=\mathfrak m B_{\mathfrak n}$. This means that there is an element $u\in B,u\not \in \mathfrak n$ such that 
\begin{equation}\label{eq:obsProof1}
    u\mathfrak n\subset \mathfrak mB.    
\end{equation}

On the other hand, by assumption, $(f_\C)^{-1}(f(X^\ns))=X^\ns$, and $f|_{X^\ns}$ is
injective. Hence $x$ is the only point of $X_\C$ lying in the fibre over $y=f(x)$, so the radical of $\mathfrak mB$ is $\mathfrak n$. Since $\mathfrak n$ is a maximal ideal this means that $\mathfrak m B$ is $\mathfrak n$-primary. From \eqref{eq:obsProof1} we deduce that $\mathfrak n\subset \mathfrak m B$, which finishes the proof of the claim.
\end{proof}

\end{document}
